\chapter{Predicción en Tres Capas}
\label{cap:prediccion}

Este capítulo combina la localización (capítulo «\nameref{cap:localizacion}») y el modelo de tiempos
(capítulo «\nameref{cap:modelo}») en la predicción que ve el usuario. La idea organizadora es una jerarquía de fuentes,
cada una con una regla explícita para abstenerse:

\begin{description}
  \item[Capa 1 · GPS.] La posición proyectada sobre el recorrido. Es la fuente más precisa y la que permite afirmar un
        adelanto o un retraso.
  \item[Capa 2 · Parada registrada.] Cuando el GPS no es utilizable pero el feed dice cuál es la próxima parada y a qué
        hora se tomó el dato, se supone al autobús a mitad del tramo que lleva a esa parada.
  \item[Capa 3 · Horario corregido.] Cuando no hay ningún dato del vehículo, se recalcula el horario desde la cabecera
        con los tiempos aprendidos.
\end{description}

\noindent Si ninguna capa puede sostener una estimación, se muestra el horario publicado, marcado como tal. Ninguna
capa inventa: cada una se usa solo en las condiciones en que se ha medido que mejora a la siguiente.

\section{Entradas, Salidas y Estados}
\label{sec:estados}

La unidad de trabajo es una \emph{fila}: una salida concreta de una parada. Una fila identifica un viaje $J$ (su
línea, número de viaje y salida de cabecera $d_1$) y una parada de ese viaje, la parada de \emph{destino} de la
predicción, con su orden $q$ y su hora publicada $a_\dest$. Para cada fila el motor devuelve:

\begin{itemize}
  \item un \textbf{estado}, de la Tabla~\ref{tab:estados};
  \item una \textbf{llegada estimada} $\ETA$, en segundos desde medianoche;
  \item un \textbf{retraso} $\delta$ en segundos, que puede ser \emph{nulo} (no se afirma retraso alguno);
  \item el vehículo y su parada actual, si se conocen.
\end{itemize}

\begin{table}[ht]
\centering
\caption{Estados de una Fila y lo que Muestra la Aplicación}\label{tab:estados}
\begin{tabularx}{\textwidth}{@{}l>{\raggedright\arraybackslash}X>{\raggedright\arraybackslash}X@{}}
\toprule
Estado & Significado & Texto en la aplicación\\
\midrule
\texttt{live} & Capa 1: el GPS sitúa al autobús en el recorrido. & «En hora», «Retraso $+N$m» o «Adelantado $-N$m»; distintivo «En ruta», «Llegando» o «En la parada» (también mientras espera en la cabecera). \\
\texttt{estimated} & Capa 2 o capa 3: estimación sin seguimiento en directo; retraso $0$. & «Estimado con lo que tardan de verdad los buses»; distintivo «Estimado». \\
\texttt{departed} & El autobús ya pasó por la parada. & «Ya salió»; distintivo «Saliendo». \\
\texttt{likely\_departed} & Sin datos del autobús, el horario corregido indica que ya debió de pasar; o, con GPS, la hora estimada de llegada pasó hace más de 3~min sin que el feed ni el GPS confirmen el paso. & «Sin datos del bus: puede que ya haya pasado»; distintivo «Sin datos del bus»; se ordena al final. \\
\texttt{scheduled} & Ninguna capa puede sostener una estimación. & «Horario, sin seguimiento en vivo»; distintivo «Solo horario». Si el feed trae autobuses con hora de hoy y todavía no es la hora del viaje, la ficha dice que aún no tiene autobús asignado. \\
\bottomrule
\end{tabularx}
\end{table}

La fila se empareja con una entrada del feed como se explicó en la sección «\nameref{sec:emparejamiento}»: primero de forma
exacta por los alias de viaje y, si no, por línea, número de viaje y salida de cabecera con una tolerancia de 180~s.
Si no hay entrada, la fila se resuelve sin datos del vehículo: con el registro de pasos (sección
«\nameref{sec:registro-pasos}») o, si este no la da por pasada, con la capa~3.

En todo el capítulo, $D$ es la fecha de servicio del viaje (de madrugada, la de ayer para los viajes que pasan de las 24:00; sección «\nameref{sec:mas-de-24h}»), $\ell$ la línea, $n$ el orden de la próxima parada que declara el
feed y $E(\cdot)$ el tiempo esperado de la definición~\ref{def:tiempo-esperado} aplicada al viaje $J$.

\section{Capa 1: Llegada Desde el GPS}
\label{sec:capa1}

Si existe la posición $(\Tpos,\tgps)$ (sección «\nameref{sec:condiciones-gps}»), la franja y el factor se toman en el
instante de la posición,
\[
b=\franja(D,\tgps),\qquad k=k(\ell,D,\tgps),
\]
y la llegada a la parada de destino es
\begin{equation}
\label{eq:eta-capa1}
\boxed{\;\ETA=\redondeo\big(\tgps+E(J,\Tpos,a_\dest,b,k)\big),\qquad \delta=\ETA-a_\dest\qquad(a_\dest\ge\Tpos).\;}
\end{equation}
Si $a_\dest<\Tpos$, la parada ya ha quedado atrás según el horario y
\begin{equation}
\label{eq:paso-capa1}
\ETA=\tgps-(\Tpos-a_\dest),\qquad\delta=\delta_{\mathrm{pos}}=\tgps-\Tpos,
\end{equation}
que es la hora a la que se estima que pasó: la distancia programada entre la parada y el autobús, contada hacia atrás
desde la hora de la posición.

La llegada parte de $\tgps$, no de $\tnow$. Cada posición llega al motor un minuto o más después de tomarse (66~s
de mediana y 92~s en el percentil 90 el 5 de noviembre de 2026; sección «\nameref{sec:llegando}»); sumar el tiempo
restante a $\tnow$ retrasaría todas las llegadas estimadas en lo que tarda en recibirse el dato. Partiendo de $\tgps$, la antigüedad
del dato no sesga la predicción: solo la hace algo menos precisa. Una consecuencia es que $\ETA$ puede ser anterior a $\tnow$ si el autobús estaba muy cerca
de la parada cuando se tomó la posición.

\section{Capa 1 en la Cabecera}
\label{sec:cabecera-eta}

Con la posición en la cabecera~\eqref{eq:cabecera}, la ecuación~\eqref{eq:eta-capa1} daría $d_1+E(J,a_1,a_\dest,b,k)$,
lo que supone que el autobús sale exactamente a su hora. Los autobuses no salen así, y la medida lo confirmó: con esa
suposición, las llegadas a 10--20~min vista se adelantaban 1~min de media y los aciertos dentro de $\pm2$~min bajaban
del 67{,}0 al 64{,}0\,\% (sección «\nameref{sec:eval-trazados}»). La llegada es, por tanto, la de la capa 3,
\begin{equation}
\label{eq:eta-cabecera}
\ETA=\max\big(C(J,a_\dest),\ \tnow\big),\qquad \delta=\ETA-a_\dest\qquad(a_\dest>d_1),
\end{equation}
con $C$ el horario corregido~\eqref{eq:capa3}, cuyos 60~s iniciales y cuya corrección por viaje
$\Delta_v$~\eqref{eq:correccion-viaje} recogen lo tarde que salen de verdad los autobuses. La fila es «en directo»
porque se sabe que el autobús está allí; si no hay tiempos aprendidos, se usa la ecuación~\eqref{eq:eta-capa1}. La
parada de cabecera ($a_\dest\le d_1$) tiene como llegada la hora de salida, $d_1$.

\section{Capa 2: Llegada Desde la Parada Registrada}
\label{sec:capa2}

Muchas entradas del feed traen la próxima parada y la hora del registro, pero no una posición utilizable: sin
coordenadas, con un GPS fuera del pasillo o con una hora de posición que no es de hoy. La capa 2 aprovecha lo que sí
es un dato real: el autobús ha dejado atrás la parada $p_{n-1}$ y va hacia $p_n$, y eso se sabía a la hora $\tref$.

\subsection{El Registro}

El \emph{registro} de una entrada existe solo si se cumplen todas estas condiciones:

\begin{enumerate}
  \item la entrada tiene viaje y orden $n$, y el par (viaje, vehículo) no está marcado como no fiable (sección
        «\nameref{sec:desconfianza}»);
  \item $n\ge2$: hay una parada anterior;
  \item hay hora de referencia $\tref$ ($\tgps$ si existe y, si no, $t_{\mathrm{rec}}$) y es reciente:
        $\tnow-\tref\le600$~s;
  \item el viaje ha empezado: $\tref\ge d_1-60$~s;
  \item el viaje tiene paradas de orden $n-1$ y $n$;
  \item el tramo es corto: $a_n-a_{n-1}\le1\,200$~s (\texttt{REGISTRATION\_MAX\_SEGMENT\_SECONDS}).
\end{enumerate}

\noindent Entonces, con $b=\franja(D,\tref)$ y $k=k(\ell,D,\tref)$,
\begin{align}
\label{eq:registro}
\theta&=\redondeo\Big(\tfrac12\,E(J,a_{n-1},a_n,b,k)\Big) &&\text{(lo que le falta para llegar a $p_n$)},\\
\delta_{\mathrm{reg}}&=\tref+\theta-a_n &&\text{(retraso estimado en $p_n$)},\\
g&=a_n-a_{n-1} &&\text{(duración programada del tramo)}.
\end{align}

\noindent Si el tramo de $p_{n-1}$ a $p_n$ tiene duración programada nula ($a_{n-1}=a_n$, frecuente cuando dos paradas
comparten minuto), $E(J,a_{n-1},a_n,b,k)=0$ por la regla de los tramos de duración nula (sección
«\nameref{sec:tiempo-esperado}»), aunque ese tramo tenga tiempo aprendido. Entonces $\theta=0$ y $g=0$: el autobús se
supone ya en $p_n$, y el contraste de la sección «\nameref{sec:desconfianza}» se hace con
$\delta_{\mathrm{reg}}=\tref-a_n$.

\subsubsection{Por Qué la Mitad} El feed no dice en qué punto del tramo está el autobús. Si el instante de la observación
es independiente de la posición, el autobús está en cualquier punto del tramo con la misma probabilidad, y el valor
esperado del tiempo que le falta es la mitad del tiempo del tramo. Se midieron repartos de 0,35 a 0,65 y el óptimo fue
0,5 (sección «\nameref{sec:descartadas}»).

\subsubsection{Por Qué Tramos de Hasta 20 Minutos} En tramos largos (salidas de cabecera, regulaciones) la incertidumbre
de «a mitad del tramo» es de muchos minutos. Contrastada con el GPS en 201 autobuses, la estimación de la parada
registrada se desvía 0,7~min de mediana y 1,7~min en el percentil 90 en tramos de hasta 5~min, pero 10~min de mediana
en tramos de más de 10~min. Con el umbral de 20~min, la capa 2 mejora al horario en todas las longitudes de tramo en que
se usa; por encima de ese umbral, no.

\subsection{La Llegada}

Para una parada de destino con orden $q\ge n$,
\begin{equation}
\label{eq:eta-capa2}
\boxed{\;\ETA=\redondeo\big(\tref+\theta+E(J,a_n,a_\dest,b,k)\big)\;}
\end{equation}
donde el último término es cero si $a_\dest\le a_n$. Para $q<n$, el registro no da llegada (la parada ya ha quedado
atrás; véase la sección «\nameref{sec:paso}»).

\section{Contraste Entre el GPS y la Parada Registrada}
\label{sec:desconfianza}

Cuando una entrada tiene a la vez posición y registro, las dos fuentes deberían coincidir. Si no coinciden, una de
ellas está equivocada (un GPS que pertenece a otro vehículo, un orden de parada que no se actualizó) y no
hay forma de saber cuál. La regla es que el par (viaje, vehículo) deja de ser fiable si

\begin{equation}
\label{eq:contraste}
g\le360\ \text{s}\qquad\text{y}\qquad\big|\delta_{\mathrm{pos}}-\delta_{\mathrm{reg}}\big|>360\ \text{s}.
\end{equation}
Los dos umbrales son las constantes \texttt{REGISTRATION\_CROSSCHECK\_MAX\_SEGMENT\_SECONDS} y
\texttt{MAX\_GPS\_VS\_STOP\_DISAGREEMENT\_SECONDS}. Un par no
fiable no tiene posición ni registro durante el resto de la petición, de modo que sus filas pasan a la capa 3. El
contraste solo se hace en tramos de hasta 6~minutos, porque en los largos el registro es demasiado impreciso para
contradecir al GPS.

Para que el resultado sea reproducible, la posición (con su contraste) debe evaluarse antes que cualquier uso del
registro. En el código, cada consulta del estado de una fila empieza por la posición; el registro que se usa dentro del
contraste se calcula antes de marcar el par y, una vez marcado, todas las consultas posteriores lo tratan como
inexistente.

\begin{ejemplo}
En el viaje sintético, un autobús declara $n=3$ (va hacia $S2$, a 1\,000~m), pero su GPS está a 2\,800~m, entre $S3$ y
$S4$, a las $29\,400$~s. La proyección elige el candidato «más allá» (a 11~m del trazado, frente a 1\,800~m del mejor
«esperado») y da $\Tpos=29\,361$, luego $\delta_{\mathrm{pos}}=39$~s. El registro da $\theta=70$ y
$\delta_{\mathrm{reg}}=29\,400+70-29\,010=460$~s. Como $g=120\le360$ y $|39-460|=421>360$, el par se marca como no
fiable y todas las filas de ese viaje se resuelven sin datos del vehículo.
\end{ejemplo}

\section{Capa 3: Horario Corregido}
\label{sec:capa3}

Sin ningún dato fiable del vehículo, la mejor estimación de su llegada es la de un autobús que sale de la cabecera poco
después de su hora y recorre cada tramo en el tiempo aprendido. Si la red tiene tabla de tiempos por tramo, el viaje
tiene paradas y $a_\dest>d_1$, se toman la franja y el factor a mitad de camino,
\[
t_{\mathrm{mid}}=\Big\lfloor\frac{d_1+a_\dest}{2}\Big\rfloor,\qquad b=\franja(D',t_{\mathrm{mid}}),\qquad
k=k(\ell,D',t_{\mathrm{mid}}),
\]
donde $D'$ es la fecha de servicio $D$, salvo en un festivo entre semana, en que es el domingo de esa semana. Las
franjas solo dependen del día de la semana, así que en un festivo se usan las de fin de semana. El \emph{horario
corregido} es
\begin{equation}
\label{eq:capa3}
\boxed{\;C=d_1+60+\redondeo\big(E(J,\,d_1-1,\,a_\dest,\,b,\,k)\big)+\redondeo(\Delta_v).\;}
\end{equation}
Los $60$~s (\texttt{SCHEDULE\_ORIGIN\_OFFSET\_SECONDS}) recogen que los autobuses salen de la cabecera algo después de
la hora publicada; el resumen de cada captura vigila que este sesgo siga siendo el medido. El inicio en $d_1-1$ incluye
los tramos de duración nula del principio del viaje (sección «\nameref{sec:tiempo-esperado}»). El último término,
$\Delta_v$, es la corrección por viaje de la sección «\nameref{sec:correccion-viaje}»; sin tabla de correcciones vale 0.

La decisión es:

\begin{enumerate}
  \item Si $C\ge\tnow$: estado \texttt{estimated}, $\ETA=C$, $\delta=0$. El autobús todavía no debería haber pasado y
        $C$ es una estimación mejor que el horario publicado.
  \item Si no, y además
        \begin{equation}
        \label{eq:probablemente}
        C<\tnow-600,\qquad a_\dest-C\ge600,\qquad \lambda\ge0{,}8
        \end{equation}
        (constantes \texttt{WRONG\_SCHEDULE\_PASSED\_SECONDS}, \texttt{WRONG\_SCHEDULE\_GAP\_SECONDS} y
        \texttt{WRONG\_SCHEDULE\_LEARNED\_SHARE}): estado \texttt{likely\_departed}, $\ETA=a_\dest$, $\delta=0$.
  \item En otro caso, estado \texttt{scheduled}, $\ETA=a_\dest$, $\delta=0$.
\end{enumerate}

\noindent La \emph{parte aprendida} $\lambda$ del camino es la fracción de su duración programada que recorre tramos
aprendidos en la franja $b$:
\begin{equation}
\label{eq:parte-aprendida}
\lambda=\frac{\sum_{i\in I(d_1-1,\,a_\dest)}\mathbf{1}\big[\tau(\kappa_i,b)\text{ definido}\big]\,h_i}{\sum_{i\in I(d_1-1,\,a_\dest)}h_i},
\qquad\lambda=0\ \text{si el denominador es }0,
\end{equation}
con el mismo conjunto de tramos $I$ de la definición~\ref{def:tiempo-esperado} y las duraciones programadas completas
$h_i$ (sin los pesos $w_i$).

Las tres condiciones de~\eqref{eq:probablemente}, que llevan al estado \texttt{likely\_departed}, identifican un caso
concreto y frecuente: un horario publicado mucho más lento que la realidad (típicamente, un viaje con un tramo largo de
autopista), en el que el horario corregido dice que el autobús pasó hace más de diez minutos, la hora publicada queda al
menos diez minutos después de esa hora corregida (a menudo, todavía no ha llegado), y la corrección se apoya en tramos medidos en al menos el 80\,\% del camino. Mostrar la
hora publicada como si el autobús fuera a venir haría que alguien esperase un autobús que ya se fue. No se oculta
porque no hay ningún dato del vehículo que lo confirme: se muestra al final de la lista, con su hora publicada y el
texto «Sin datos del bus: puede que ya haya pasado», y las alertas de llegada lo tratan como pasado. El texto dice lo que
se sabe: en las capturas del 3 de noviembre de 2026, el 97{,}8\,\% de los viajes con este estado no apareció en el feed
en todo el día (sección «\nameref{sec:fiabilidad-estados}»).

\subsection{Corrección por Viaje}
\label{sec:correccion-viaje}

Sin la corrección, el horario corregido llegaba de media 131~s antes que el autobús (sección
«\nameref{sec:eval-correccion}»). Hay dos razones. La suma de las medianas de los tramos queda por debajo del tiempo
medio del recorrido, porque el tiempo de cada tramo tiene una cola larga hacia arriba. Además, algunos viajes salen tarde
o se retrasan en el mismo sitio casi todos los días. La corrección $\Delta_v$ recoge cuánto se aparta habitualmente cada
viaje del horario corregido.

Un viaje se reconoce por la clave $v=(\ell,\,p_1,\,p_N,\,d_1,\,N)$: su línea, su primera y su última parada, su
salida de cabecera y su número de paradas. A diferencia del identificador del GTFS, esta clave no cambia cuando el
operador publica un horario nuevo con los mismos viajes. Como en el factor de ritmo (sección «\nameref{sec:k}»), la
línea, tanto en la clave como en la tabla por línea, se compara sin ceros a la izquierda. El tipo de día $\theta$ es
laborable o fin de semana según $D'$, de modo que un festivo entre semana cuenta como fin de semana. La parada de destino, en la posición $i$ del
recorrido (0 es la cabecera), cae en el tercio
\begin{equation}
\label{eq:tercio}
p=\min\big(2,\ \lfloor 3i/N\rfloor\big).
\end{equation}
La corrección se busca de lo más fino a lo más general:
\begin{equation}
\label{eq:correccion-viaje}
\Delta_v=\begin{cases}
\Delta(v,\theta,p) & \text{si el viaje tiene valor para ese tipo de día},\\
l(\ell,b_1) & \text{si no, y su línea lo tiene en la franja } b_1=\franja(D',d_1),\\
g(b_1) & \text{si no, y la franja lo tiene},\\
0 & \text{en otro caso.}
\end{cases}
\end{equation}
Las tres tablas se aprenden cada día (sección «\nameref{sec:recalibracion}»). Un día es festivo entre semana si el
horario hace circular ese día, por excepción, algún calendario de solo domingo (máscara de días $64$), que es como
publica Bizkaibus sus festivos. La corrección solo afecta a la capa 3: las capas 1 y 2 parten de donde está el autobús
y ya recogen lo que le pasa ese día.

\subsection{Sin Conexión}
\label{sec:sin-conexion}

La aplicación puede guardar en el móvil el horario completo de cada red para usarlo sin cobertura. En Bizkaibus, cada
viaje lleva además su horario corregido~\eqref{eq:capa3}, ya con $\Delta_v$, calculado una vez para un laborable y otra
para un fin de semana. El móvil usa el de fin de semana en sábado, en domingo y en los festivos, que reconoce con la
misma regla de los calendarios de solo domingo. Aplica las reglas de estado del servidor, con la parte aprendida
$\lambda$ calculada sobre el recorrido entero: \texttt{estimated} si $C\ge\tnow$ y \texttt{likely\_departed} con las
condiciones~\eqref{eq:probablemente}. En 1\,200 paradas al azar de un laborable, un sábado y un festivo, las horas del
móvil y del servidor difieren como mucho 6~s, por el orden de los redondeos y por cómo se cuenta el primer tramo
cuando la salida de cabecera va unos segundos después de su llegada.

\section{Paradas Ya Pasadas}
\label{sec:paso}

Antes de calcular una llegada, se comprueba si el autobús ya pasó por la parada de destino. La hora de paso
$t_{\mathrm{paso}}$ se calcula así:

\begin{itemize}
  \item Con posición en la cabecera (sección «\nameref{sec:cabecera}»): el autobús aún no ha pasado por ninguna
        parada.
  \item Con otra posición: si $a_\dest<\Tpos$ y además el feed ya ha dejado atrás esa parada ($q<n$) o el GPS está a más
        de 150~m de ella (\texttt{PASSED\_DISTANCE\_METERS}), $t_{\mathrm{paso}}=\tgps-(\Tpos-a_\dest)$
        (ecuación~\eqref{eq:paso-capa1}); si no, pero $q<n$ (el feed ya ha dejado atrás esa parada aunque la proyección
        ponga al autobús antes de ella), $t_{\mathrm{paso}}=\tgps$; si no, no ha pasado.
  \item Sin posición, con registro y $q<n$: $t_{\mathrm{paso}}=\tref-\max(0,\ a_{n-1}-a_\dest)$.
\end{itemize}

\noindent Si hay hora de paso, la fila sale como \texttt{departed} con $\ETA=\min(t_{\mathrm{paso}},\tnow)$ y
$\delta=0$.

\section{Llegando sin Confirmar el Paso}
\label{sec:llegando}

El feed de Bizkaibus publica una versión cada 64~s de mediana (de 58 a 72~s) y cada posición llega 66~s después de
tomarse (mediana; 92~s en el percentil 90), según las capturas del 5 de noviembre de 2026 de 7:00 a 9:00. La posición
que se usa es, por tanto, de hace un minuto o más. Si el autobús va algo más lento de lo esperado, la llegada calculada
desde esa posición con la ecuación~\eqref{eq:eta-capa1} puede quedar ya en el pasado sin que el autobús haya llegado a
la parada. Antes, la fila salía en directo con esa hora ya pasada y la lista la retiraba al quedar atrás (sección
«\nameref{sec:retraso-mostrado}»): en esas dos horas, el autobús desaparecía de la lista antes de llegar en 170 de las
5\,041 llegadas a parada que el GPS ve a menos de 40~m (el 3{,}4\,\%).

Con posición y sin hora de paso, si $\ETA<\tnow$, la fila sale con $\ETA=\tnow$ (en la aplicación, «0~min») y un
retraso de al menos $\tnow-a_\dest$, hasta que el feed o el GPS confirmen el paso. Si $\tnow-\ETA$ pasa de 180~s
(\texttt{ARRIVING\_MAX\_SECONDS}) sin confirmación, la fila sale como \texttt{likely\_departed} («puede que ya haya
pasado») con $\ETA=\tnow$, en vez de desaparecer. Con las mismas capturas, el autobús desaparece antes de llegar en
22 de las 5\,041 llegadas (el 0{,}4\,\%; recorridos que pasan dos veces cerca de la parada o un feed que salta
paradas) y, a cambio, tarda algo más en dejar de mostrarse después de irse: 104~s de mediana y 152~s en el percentil
90, frente a 61 y 98~s. La ficha del viaje usa la misma regla de paso y la misma llegada. El estado \texttt{likely\_departed} tiene
así dos usos: en la capa 3 lleva la hora publicada y retraso 0; aquí, la hora actual y el retraso que da el GPS. Quien
use la API no debe tratar igual su llegada en los dos casos.

\section{Registro de Pasos}
\label{sec:registro-pasos}

Un viaje que ya terminó, o que ha dejado de aparecer en el feed, no tiene entrada, y sus filas irían a la capa 3 como
si el autobús no hubiera salido. Para evitarlo, un proceso externo guarda cada minuto, para cada referencia de viaje
del feed, el último orden de parada visto y su hora:
\[
\text{pasos}[D][\rho]=\big(n,\ t\big),\qquad t=t_{\mathrm{rec}}\ \text{si existe; si no, }\tgps,
\]
donde $t$ es la hora del registro en que el feed dio ese orden de parada (la capa 2 y la memoria del feed prefieren, en
cambio, $\tgps$) y $D$ es el día del viaje, no el del reloj: pasada la medianoche, un autobús de un viaje de ayer se
apunta en los pasos de ayer, con $t$ en las horas de su horario (sección «\nameref{sec:mas-de-24h}»). Hasta el 7 de
noviembre de 2026 se apuntaba en el día del reloj: el autobús de las 23:30 de ayer, visto a las 00:30, quedaba en los
pasos de hoy, y las paradas del viaje de las 23:30 de hoy, que tiene la misma referencia, salían esa noche como ya
pasadas hasta su salida. Los pasos de cada día caducan 36~horas después de terminar ese día. Al responder, y solo para
los viajes de las filas que no tienen entrada en el feed, se leen los pasos de todas sus referencias (un viaje
fusionado puede tener varias; sección «\nameref{sec:alias}») y se conserva el de mayor orden. Si la fila tiene $q<n$,
sale como \texttt{departed} con $\ETA=\min(t,\tnow)$. Esta comprobación se hace antes de recurrir a la capa 3.

\section{Autobuses que Faltan en una Lectura del Feed}
\label{sec:memoria-feed}

El feed deja fuera a veces, durante una lectura o unas pocas, a un autobús que sigue en su viaje. En las capturas del 2
al 4 de noviembre de 2026 pasa unas 700 veces al día con una sola lectura y otras 250 a 400 con dos o tres. Sin otra
medida, en esas lecturas las filas del viaje no tienen entrada y se resuelven con la capa 3, o salen como
\texttt{scheduled} («Solo horario») si el horario corregido ya ha pasado, aunque el autobús se viera un minuto antes.

Por eso, al leer el feed, se le añade la última entrada de cada autobús que falta en la lectura nueva si se cumplen
las tres condiciones siguientes:
\begin{itemize}
  \item el dato tiene como mucho $H=180$~s de antigüedad (\texttt{HOLD\_SECONDS}): $0\le\tnow-\tgps\le H$, con
        $t_{\mathrm{rec}}$ si no hay posición;
  \item el autobús no aparece en la lectura nueva con otro viaje;
  \item su viaje no lo lleva ya otro autobús en la lectura nueva.
\end{itemize}
La entrada se usa tal cual, y la capa~1 parte de su $\tgps$~\eqref{eq:eta-capa1}. Como ya se explicó, la antigüedad del
dato no sesga la llegada: solo la hace algo menos precisa. La memoria es un fichero temporal por red en el servidor;
en un servidor recién arrancado no hay nada guardado y se usa solo la lectura nueva.

\begin{table}[ht]
\centering
\caption{Llegadas de Autobuses que Faltaban en el Feed, con y sin Guardar su Último Dato}\label{tab:memoria-feed}
\begin{tabular}{@{}lrrrr@{}}
\toprule
Llegada en & $n$ & Sin guardar & Guardando 180~s & Mediana del error\\
\midrule
0--5 min & 9\,572 & 54{,}5\,\% & 77{,}5\,\% & de 1{,}8 a 1{,}0 min\\
5--10 min & 8\,945 & 51{,}6\,\% & 65{,}2\,\% & \\
10--20 min & 8\,676 & 37{,}3\,\% & 42{,}8\,\% & \\
\bottomrule
\end{tabular}
\nota{Aciertos dentro de $\pm2$~min contra el paso real, con las capturas del 2 al 4 de noviembre de 2026: cada
lectura en la que falta un autobús visto en los 180~s anteriores y que vuelve después al mismo viaje.}
\end{table}

La Tabla~\ref{tab:memoria-feed} recoge el resultado. Guardar el dato mejora los aciertos en los tres días y en los tres
horizontes; con 60, 120 y 300~s, la mejora a 0--5~min es del 55{,}7 al 80{,}9, del 54{,}8 al 79{,}6 y del 54{,}3 al
76{,}0\,\% (cada umbral mide sus propias lecturas, las de autobuses vistos en los $H$ segundos anteriores, y por eso
cambia también la cifra sin guardar). Cuanto más antiguo es el dato, algo menos acierta, y 180~s cubre la mayoría de
los cortes. El coste es el autobús que desaparece de verdad (acaba o se pierde el viaje): sus paradas siguen «en
directo» hasta que su último dato cumple $H$ segundos, cuando sin guardarlo habrían salido con la capa~3, que tampoco
se cumple.

\section{Viajes que Pasan de las 24:00}
\label{sec:mas-de-24h}

El horario cuenta las horas desde la medianoche del día del viaje: un autobús del viernes que pasa a las 02:10 del
sábado figura a las 26:10 ($94\,200$~s). En el horario de Bizkaibus del 5 de noviembre de 2026, 52 viajes pasan de las 24:00, y 43 de ellos salen ya pasada esa
hora: 6 circulan todas las noches, hasta las 26:05; 33, solo las noches del viernes o del sábado, y llegan hasta las
32:10, las 08:10 del día siguiente; y 4, solo en fechas concretas. El feed, en cambio, da $\tgps$ y $t_{\mathrm{rec}}$ en horas del día en curso, y la referencia
de un viaje que circula todos los días es la misma cada día. Con un solo emparejador, el autobús de las 23:30 de ayer,
visto a las 00:30, se emparejaría con el viaje de las 23:30 de hoy, y no con el suyo.

Por eso, cada entrada del feed se asigna al día de su viaje. Una entrada es de un viaje de ayer si la salida de cabecera
de su referencia queda más de 12~h por delante de la hora actual (\texttt{RealtimeMatcher::daysAgo}):
\begin{equation}
\label{eq:dia-viaje}
d_0^{\mathrm{feed}}-\tnow>43\,200\ \text{s}.
\end{equation}
A las 00:30 ($\tnow=1\,800$) son de ayer el viaje que salió a las 23:30 ($d_0^{\mathrm{feed}}=84\,600$) y el de las
24:10 ($87\,000$). La regla supone que el feed solo trae viajes en marcha o a punto de salir, y que la referencia da la
salida en horas del horario, como la de Bizkaibus~\eqref{eq:regex}.

Las entradas de ayer van a un emparejador propio en el que todo está en las horas del día del viaje: $\tgps$,
$t_{\mathrm{rec}}$ y $\tnow$ se cuentan con $86\,400$~s más (las 02:10 son las 26:10), la fecha de servicio $D$ es la de ayer y
el registro de pasos, el de ayer (\texttt{RealtimeMatcher::vehiclesOfDay}). El cálculo es el de las secciones
anteriores, sin cambios: la franja de las 26:10 de un viernes, por ejemplo, es la \texttt{laborable.madrugada} del
viernes (definición~\ref{def:franja}). El emparejador de hoy solo recibe las entradas de hoy, y una entrada sin ninguna
hora de hoy no va a ninguno (sección «\nameref{sec:feed}»).

La aplicación pasa al emparejador de ayer las salidas de ayer que pasan de las 24:00, con sus horas del horario, y
devuelve sus llegadas en horas de hoy, $\ETA-86\,400$ (sección «\nameref{sec:ultima-salida}»). De madrugada, la ficha de
un viaje muestra el viaje de ayer mientras este no haya terminado, es decir, hasta 10~min después de su última parada
(\texttt{LATE\_GRACE\_SECONDS}); en la red de autobús, solo si el feed no trae ningún autobús en el viaje de hoy. En el
mapa de una línea, cada autobús se sitúa con el emparejador de su día.

En la práctica, de 23:58 a 05:00 el feed de Bizkaibus no se actualiza (sección «\nameref{sec:feed}»), así que los
viajes nocturnos salen con la capa~3 o con el horario; el mapa no dice que siga a ningún autobús, y la ficha no dice
que a un viaje le falte autobús. Antes, el mapa de la línea 3136 mostraba toda la noche seis autobuses vistos a las
23:58. El emparejador de ayer sirve para los feeds que siguen actualizándose de noche y, en Bizkaibus, para los 11
viajes de los horarios del viernes y del sábado que acaban después de las 29:00 (las 05:00 del día siguiente), en las
líneas 3923 y 3932. De día, ninguna entrada es de ayer: con el código anterior y el nuevo sobre el mismo feed, 28 paradas, 256 fichas y 56 mapas de línea
dan las mismas respuestas.

\section{La Decisión Completa}
\label{sec:decision}

El algoritmo de la Figura~\ref{alg:fila} reúne las reglas anteriores en el orden exacto en que se aplican. Cualquier error
inesperado durante el cálculo de una fila (un dato corrupto, una base de datos inaccesible) se captura y la fila sale
con la capa 3 o con el horario publicado: una fila nunca impide responder a las demás.

\begin{figure}[ht]
\caption{Estado y Llegada de una Fila (Motor de Autobús)}\label{alg:fila}
\begin{algorithmic}[1]
\Require fila $(J,q,a_\dest)$, entrada del feed $e$ (o ninguna), $\tnow$
\Ensure (estado, $\ETA$, $\delta$)
\State $P\gets$ posición de $e$ en $J$ con su contraste (capítulo «\nameref{cap:localizacion}» y sección «\nameref{sec:desconfianza}»)
\State $G\gets$ registro de $e$ en $J$ (sección «\nameref{sec:capa2}»)
\If{$e$ no existe \textbf{o} ($P$ no existe \textbf{y} $G$ no existe)}
  \State \Return \Call{SinDatos}{$J,q,a_\dest$}
\EndIf
\If{hay $t_{\mathrm{paso}}$ (sección «\nameref{sec:paso}»)}
  \State \Return $(\texttt{departed},\ \min(t_{\mathrm{paso}},\tnow),\ 0)$
\EndIf
\If{$P$ no existe \textbf{y} ($q$ desconocido \textbf{o} la capa 2 no da llegada para $q$)}
  \State \Return \Call{SinDatos}{$J,q,a_\dest$}
\EndIf
\If{$P$ existe}
  \State $(\ETA,\delta)\gets$ capa 1, ecuaciones~\eqref{eq:eta-capa1}, \eqref{eq:paso-capa1} o, en la cabecera,~\eqref{eq:eta-cabecera}; $\pi\gets$ verdadero
\Else
  \State $\ETA\gets$ capa 2, ecuación~\eqref{eq:eta-capa2}; $\delta\gets\ETA-a_\dest$; $\pi\gets$ falso
\EndIf
\If{$|\delta|>2\,700$} $\delta\gets$ nulo \EndIf \Comment{retraso no creíble: se conserva la llegada}
\State $\alpha\gets(\ETA<\tnow-120)$ \Comment{ya debería haber pasado}
\If{$\alpha$ \textbf{y no} $\pi$} \State \Return \Call{SinDatos}{$J,q,a_\dest$} \EndIf
\If{\textbf{no} $\pi$} \State \Return $(\texttt{estimated},\ \ETA,\ 0)$ \EndIf
\If{$\ETA<\tnow$} \Comment{llegando, sin confirmar el paso}
  \If{$\delta$ no es nulo} $\delta\gets\max(\delta,\ \tnow-a_\dest)$ \EndIf
  \If{$\tnow-\ETA>180$} \State \Return $(\texttt{likely\_departed},\ \tnow,\ \delta)$ \EndIf
  \State \Return $(\texttt{live},\ \tnow,\ \delta)$
\EndIf
\State \Return $(\texttt{live},\ \ETA,\ \delta)$
\Statex
\Function{SinDatos}{$J,q,a_\dest$}
  \If{$e$ no existe \textbf{y} el registro de pasos tiene $(n,t)$ para $J$ con $q<n$} \State \Return $(\texttt{departed},\min(t,\tnow),0)$ \EndIf
  \State \Return capa 3 (sección «\nameref{sec:capa3}»): \texttt{estimated}, \texttt{likely\_departed} o \texttt{scheduled}
\EndFunction
\end{algorithmic}
\end{figure}

\subsection{El Tope de 45 Minutos} Un desfase de más de 45~minutos, de retraso o de adelanto ($|\delta|>2\,700$~s,
\texttt{MAX\_PLAUSIBLE\_DELAY\_SECONDS}), casi siempre indica un emparejamiento erróneo, no un retraso real: en los datos, los retrasos reales llegan a $+37$~min y
los viajes mal asignados dan valores como $+111$~min. Por encima del tope el retraso no se afirma ($\delta$ nulo), pero
la llegada se conserva, porque sigue saliendo de los datos del propio autobús y es útil. Se probó a ocultar también
la llegada, y la precisión empeoraba (del 85,2\,\% al 84,1\,\% de aciertos dentro de $\pm2$~min).

\subsection{El Margen de 120 Segundos} Una llegada de la capa~2 hasta dos minutos en el pasado
(\texttt{STALE\_GRACE\_SECONDS}) se sigue mostrando como estimada: el autobús puede estar en la parada o saliendo de
ella. Más allá, la estimación ya no es creíble y la fila pasa a la capa~3. Con GPS no hay margen: una llegada ya
pasada sale con $\ETA=\tnow$ hasta que se confirme el paso (sección «\nameref{sec:llegando}»).

\section{Retraso y Minutos que Se Muestran}
\label{sec:retraso-mostrado}

El desfase en la posición $\delta_{\mathrm{pos}}$ no es el retraso que entiende el usuario: compara al autobús con un
horario interpolado en un punto donde no hay parada. Lo que el usuario compara es la hora publicada de \emph{su} parada
con la llegada que se le muestra. Por eso, el retraso que se muestra es siempre una diferencia entre una llegada
estimada y una hora publicada:

\begin{itemize}
  \item \textbf{En la lista de una parada}, para cada fila,
        \begin{equation}
        \label{eq:retraso-lista}
        \text{retraso}=\redondeo\!\big(\delta/60\big)\ \text{min},\qquad \text{minutos}=\redondeo\!\big((\ETA-\tnow)/60\big),
        \end{equation}
        con el retraso nulo si $\delta$ es nulo. Las filas con menos de $-2$ minutos se retiran de la lista. Las demás
        se ordenan por $(\mathbf{1}[\text{estado}=\texttt{likely\_departed}],\ \ETA)$: por llegada estimada, no por la
        hora publicada (un autobús adelantado puede llegar antes que uno anterior en el horario), y con las filas
        \texttt{likely\_departed} al final.
  \item \textbf{En la ficha de un autobús}, el retraso es el de su próxima parada: con posición, $\ETA(p_n)-a_n$
        calculado con la capa 1. Se muestra $\redondeo(\cdot/60)$ minutos, o nada si $|\delta|$ pasa de 45~minutos. Sin
        posición se calcula $\delta_{\mathrm{reg}}$ (o 0 si tampoco hay registro), pero no se afirma: la ficha se marca
        como estimada y sin retraso, igual que la lista.
\end{itemize}

\noindent La aplicación traduce el retraso $r$ en minutos a un texto: «Retraso $+r$m» si $r>0$, «Adelantado $r$m» si
$r\le-2$, con su signo («Adelantado $-2$m»), y «En hora» en otro caso. La asimetría es deliberada: un autobús que llega un minuto antes no se anuncia como
adelantado, porque la resolución del dato es de un minuto. Nótese el efecto del redondeo de PHP:
$\redondeo(-90/60)=\redondeo(-1{,}5)=-2$, de modo que un adelanto de 90~s ya se anuncia.

Con esta definición, la lista y la ficha nunca se contradicen: se comprobó con el feed real que ninguna de 232 salidas
mostraba un retraso incoherente con su llegada.

\section{Después de la Última Salida del Día}
\label{sec:ultima-salida}

La lista se construye con el calendario del día y, de madrugada, con las salidas del día anterior que pasan de las
24:00, con el calendario de ayer y sus horas pasadas a las de hoy ($a-86\,400$: las 26:10 son las 02:10), intercaladas
con las de hoy por hora de salida. Las de la red de autobús pasan por el emparejador de ayer (sección
«\nameref{sec:mas-de-24h}») y las de las demás redes, por el emparejamiento por viaje de la sección
«\nameref{sec:otras-redes}». Cuando ya no queda ninguna salida por delante (las que acaban de salir siguen en la lista
por si alguna viene con retraso), se añaden detrás las primeras del día siguiente, con su calendario y sus horas
desplazadas 24~h ($a+86\,400$), en estado \texttt{scheduled} y sin pasar por el motor: el feed de hoy no dice nada de
esos viajes. La aplicación las muestra
con la hora («mañana 07:00») en lugar de los minutos.

\section{Ejemplo Completo}
\label{sec:ejemplo-completo}

La Tabla~\ref{tab:escenarios} recoge los dieciséis escenarios de los vectores de prueba (Apéndice~\ref{ap:vectores})
sobre el viaje sintético del ejemplo~\ref{ej:posicion}, con los tramos aprendidos, las proporciones y el factor del
ejemplo~\ref{ej:tiempo-esperado}. Las horas publicadas son $28\,800$, $28\,890$, $29\,010$, $29\,310$ y $29\,430$~s
(8:00:00 a 8:10:30), salvo en el escenario~8. Todas las cifras se han obtenido con el código de producción.

\begin{table}[ht]
\centering
\caption{Estado y Llegada Estimada en Cada Parada del Viaje Sintético}\label{tab:escenarios}
\setlength{\tabcolsep}{3.5pt}
\begin{tabular}{@{}lrccccc@{}}
\toprule
Esc. & $\tnow$ & $S0$ & $S1$ & $S2$ & $S3$ & $S4$\\
\midrule
1 & 28\,980 & S 28\,800 & S 28\,890 & V 29\,020 (+10) & V 29\,220 ($-90$) & V 29\,318 ($-112$)\\
2 & 28\,980 & S 28\,800 & S 28\,890 & V 29\,020 (+10) & V 29\,290 ($-20$) & V 29\,398 ($-32$)\\
3 & 28\,980 & S 28\,860 & S 28\,950 & E 29\,020 & E 29\,220 & E 29\,318\\
4 & 28\,980 & S 28\,860 & S 28\,950 & E 29\,020 & E 29\,220 & E 29\,318\\
5 & 29\,420 & P 28\,800 & P 28\,890 & P 29\,010 & P 29\,310 & P 29\,430\\
6 & 28\,700 & P 28\,800 & E 28\,970 & E 29\,110 & E 29\,310 & E 29\,408\\
7 & 29\,600 & P 28\,800 & P 28\,890 & P 29\,010 & P 29\,310 & P 29\,430\\
8 & 31\,200 & P 28\,800 & P 28\,890 & P 29\,010 & Pp 30\,810 & Pp 30\,930\\
9 & 28\,980 & P 28\,800 & P 28\,890 & E 29\,110 & E 29\,310 & E 29\,408\\
10 & 28\,700 & P 28\,800 & E 28\,990 & E 29\,150 & E 29\,350 & E 29\,468\\
11 & 28\,700 & P 28\,800 & E 28\,930 & E 29\,110 & E 29\,302 & E 29\,424\\
12 & 28\,210 & V 28\,800 (0) & V 28\,970 (+80) & V 29\,110 (+100) & V 29\,310 (0) & V 29\,408 ($-22$)\\
13 & 28\,210 & P 28\,800 & E 28\,970 & E 29\,110 & E 29\,310 & E 29\,408\\
14 & 29\,100 & S 28\,800 & S 28\,890 & V 29\,100 (+90) & V 29\,220 ($-90$) & V 29\,318 ($-112$)\\
15 & 29\,040 & S 28\,800 & S 28\,890 & V 29\,040 (+30) & V 29\,214 ($-96$) & V 29\,311 ($-119$)\\
16 & 29\,320 & S 28\,800 & S 28\,890 & Pp 29\,320 (+310) & V 29\,320 (+10) & V 29\,320 ($-110$)\\
\bottomrule
\end{tabular}
\nota{Esc.: escenario. Llegadas en segundos. V: \texttt{live}; E: \texttt{estimated}; S: \texttt{departed}; P: \texttt{scheduled}; Pp: \texttt{likely\_departed}. Entre paréntesis, el retraso en segundos cuando se afirma. Escenarios: 1, GPS a 700~m del origen y 11~m del trazado; 2, el mismo sin trazado; 3, GPS a 60~m del trazado; 4, sin GPS, con $n=3$ y $\tref=28\,950$; 5, GPS y registro discrepan; 6 y 7, sin datos del autobús, antes de salir y con el horario corregido ya pasado, respectivamente; 8, horario publicado muy lento ($S3$ y $S4$ a $30\,810$ y $30\,930$); 9, GPS de hace 680~s; 10, como el 6, con la corrección por viaje de los vectores de prueba; 11, el 10 en un festivo; 12, el autobús espera en la primera parada 10~min antes de salir; 13, el mismo a 600~m de ella; 14, como el 1 con $\tnow=29\,100$: la llegada a $S2$ ya pasó sin confirmar el paso; 15, GPS 50~m pasada $S2$ con el feed aún en ella (no basta para darla por pasada); 16, como el 14 con $\tnow=29\,320$: más de 3~min llegando sin confirmar el paso.}
\end{table}

\subsection{Escenario 1} Es el del ejemplo~\ref{ej:posicion}: $\Tpos=\tgps=28\,950$. Las paradas $S0$ y $S1$ tienen
$a_\dest<\Tpos$ y salen como pasadas a las $28\,950-(28\,950-28\,800)=28\,800$ y $28\,890$. Para $S2$, $S3$ y $S4$ la
llegada es $\tgps$ más los tiempos esperados del ejemplo~\ref{ej:tiempo-esperado}: $28\,950+70=29\,020$,
$\redondeo(28\,950+270{,}4)=29\,220$ y $\redondeo(28\,950+367{,}72)=29\,318$. El horario publicado de $S3$ y $S4$ es
90 y 112~s más lento de lo esperado, y la aplicación anunciaría «Adelantado $-2$m» en ambas
($\redondeo(-1{,}5)=-2$ y $\redondeo(-1{,}87)=-2$).

\subsection{Escenario 2} Sin trazado, la proyección sobre la recta $S1S2$ da el mismo $\Tpos$, pero los tramos sin
aprender ya no tienen longitud y caen al factor de ritmo: las llegadas a $S3$ y $S4$ pasan a $29\,290$ y $29\,398$.

\subsection{Escenarios 3 y 4} Con el GPS a 60~m (fuera del pasillo de 50~m) o sin GPS, se usa el registro: $n=3$,
$\tref=28\,950$, $\theta=\redondeo(0{,}5\cdot140)=70$. La llegada a $S2$ es $28\,950+70=29\,020$, y a $S3$,
$\redondeo(28\,950+70+200{,}4)=29\,220$: las mismas cifras que con GPS, porque en este ejemplo el autobús estaba
justo a mitad del tramo. Las paradas pasadas salen a $\tref-\max(0,a_2-a_\dest)$: $28\,950-90=28\,860$ para $S0$ y
$28\,950$ para $S1$. Las demás salen como estimadas y sin retraso.

\subsection{Escenario 5} El contraste de la sección «\nameref{sec:desconfianza}» invalida GPS y registro. La capa 3 da para
$S2$ un horario corregido de $28\,800+60+250=29\,110<\tnow$, que no cumple las condiciones~\eqref{eq:probablemente}
de \texttt{likely\_departed}, de modo que todas las paradas salen como \texttt{scheduled}, con su hora publicada.

\subsection{Escenario 6} Sin datos y antes de la salida, el horario corregido es $28\,860+E$: $28\,970$ para $S1$
(110~s aprendidos), $29\,110$ para $S2$ (110+140), $29\,310$ para $S3$ ($250+200{,}4=450{,}4$, redondeado a 450) y
$29\,408$ para $S4$ ($450{,}4+97{,}32=547{,}72$, redondeado a 548). La cabecera $S0$ no tiene $a_\dest>d_1$ y sale como
\texttt{scheduled}.

\subsection{Escenario 8} Con todos los tramos aprendidos (110, 140, 240 y 100~s) y un horario publicado que pone
$S3$ a las $30\,810$, el horario corregido es $C=28\,860+490=29\,350$. A las $31\,200$ se cumple $C<31\,200-600$,
$a_\dest-C=1\,460\ge600$ y $\lambda=1$: estado \texttt{likely\_departed}. Para $S2$, $C=29\,110$ es posterior a su hora
publicada y la regla no se aplica.

\subsection{Escenario 9} Un dato de hace 680~s supera la antigüedad máxima de 600~s: no hay posición ni registro, y
las filas van a la capa 3 como en el escenario 6, salvo $S1$, cuyo horario corregido ($28\,970$) ya es anterior a
$\tnow=28\,980$.

\subsection{Escenarios 10 y 11} Con la tabla de corrección de los vectores de prueba (Apéndice~\ref{ap:vectores}), el
viaje tiene en laborable 20, 40 y 60~s para sus tres tercios. Las paradas $S1$ a $S4$ ocupan las posiciones 1 a 4 de
cinco y caen en los tercios 0, 1, 1 y 2~\eqref{eq:tercio}: las llegadas del escenario 6 se desplazan 20, 40, 40 y
60~s. En el festivo, $D'$ es el domingo. La tabla de tramos no tiene valores en franjas de fin de semana, así que se
usan los de cualquier hora (\texttt{all}: 100 y 150~s); el factor vale 1, los tramos sin aprender usan las proporciones
\texttt{all} y la corrección de fin de semana es de $-30$, 0 y 30~s. Para $S4$, $C=28\,860+\redondeo(100+150+192+92{,}28)+30=29\,424$.

\subsection{Escenarios 12 y 13} Diez minutos antes de salir, el autobús está a 11~m de la primera parada con $n=2$: el feed
ya da la segunda como próxima parada. Queda situado en la cabecera (sección «\nameref{sec:cabecera}»), con
$\Tpos=a_1$ y la hora de la posición igual a la de salida, y todas las paradas salen en directo, $S0$ incluida, que no
cuenta como pasada aunque $n=2$.
Las llegadas son las de la capa 3 del escenario 6, porque el autobús aún no ha recorrido nada; el retraso que se
afirma es el de esas llegadas frente al horario publicado. A 600~m de la parada (escenario 13) puede estar acabando
otro viaje: el autobús no se sitúa en la cabecera, y las filas son las del escenario 6.
